(1) 由电源输入瞬时功率 $p=i_1\dot\psi_1+i_2\dot\psi_2$，在 $L_1,L_2,M$ 恒定时整理为全导数 $d[\frac12 L_1i_1^2+\frac12 L_2i_2^2+Mi_1i_2]/dt$。取零电流时能量为零，得到 $U_B=\frac12\boldsymbol i^T\mathsf L\boldsymbol i$，$\mathsf L=\begin{pmatrix}L_1&M\\M&L_2\end{pmatrix}$；交叉项的系数不是 $2M$。
(2) 对任意非零电流严格正能需要 $L_1>0$、$\det\mathsf L=L_1L_2-M^2>0$，即 $|M|<\sqrt{L_1L_2}$。若允许理想完全耦合导致某模磁场为零，可放宽为 $|M|\le\sqrt{L_1L_2}$。常定义耦合系数 $k_c=M/\sqrt{L_1L_2}$，其绝对值不超过一；超出范围意味着能量矩阵违反被动系统物理约束。
(3) 串联后两个线圈流过同一电流，若点号所定互感磁链相助，则 $U_B=(L_1+L_2+2|M|)i^2/2$，$L_{\rm aid}=L_1+L_2+2|M|$；反接一只线圈使交叉项改号，$L_{\rm opp}=L_1+L_2-2|M|$。由正定条件 $L_{\rm opp}\ge0$，并且两接法的电感差为 $4|M|$，可作为互感实验测量式。
(4) 设电荷变量 $q_j$、$i_j=\dot q_j$，两个无损回路满足 $L_1\ddot q_1+M\ddot q_2+q_1/C_1=0$ 与 $M\ddot q_1+L_2\ddot q_2+q_2/C_2=0$。代入 $q_j=Q_je^{i\omega t}$，得到 $(C_1^{-1}-\omega^2L_1)(C_2^{-1}-\omega^2L_2)-M^2\omega^4=0$。在 $|M|<\sqrt{L_1L_2}$ 时两根 $\omega^2$ 均正，不会出现无源指数发散模。
